% !TEX root = ../main.tex
%
% Chapter 07 — Golden Sprout: Fibonacci Recurrence as a Morphism in Z[φ]
% Trinity S³AI — Flos Aureus v6.2
% Author: Dmitrii Vasilev <admin@t27.ai>
% Lane: L7  Branch: feat/phd-ch07  INV: INV-12
% Trinity anchor: φ² + φ⁻² = 3  DOI 10.5281/zenodo.19227877
%
% Local macro definitions (R6 compliance — all φ-derived constants via \phipow)
% These are local to this chapter and do not conflict with main.tex.
%
\providecommand{\phipow}[1]{%
  \ensuremath{\varphi^{#1}}%
}
%
% \coqcite{theorem}{file}{lines}{status}
\providecommand{\coqcite}[4]{%
  \begin{quote}\small
  \textbf{Coq:} \texttt{#1} in \filepath{#2},
  lines~#3 — \emph{#4}.
  \end{quote}%
}

\chapter{Golden Sprout: Fibonacci Recurrence as a Morphism in \(\mathbb{Z}[\varphi]\)}
\label{ch:golden-sprout}

\begin{figure}[H]
\centering
\makebox[\linewidth][c]{%
  \includegraphics[width=1.18\linewidth,keepaspectratio]{07-golden-sprout.png}}
\caption*{Figure~7 — Golden Sprout: the Fibonacci spiral and its algebraic skeleton
in the ring \(\mathbb{Z}[\varphi]\). Every arm count is a Fibonacci number;
every spacing ratio tends to \(\varphi = \phipow{1}\).}
\end{figure}

% ============================================================
\section{Abstract}\label{sec:ch7-abstract}
% ============================================================

This chapter establishes the Fibonacci sequence as an object with three
simultaneous identities: (I)~a recurrence defined by the \emph{growth rule}
\(F_{n+1} = F_n + F_{n-1}\); (II)~a ring morphism on the quadratic integer
ring \(\mathbb{Z}[\varphi]\), where \(\varphi = (1+\sqrt{5})/2\) is the
unique positive root of \(x^2 = x + 1\); and (III)~a sequence that lifts,
modulo the minimal polynomial of \(\varphi^4\), into a finite field of
characteristic 2 (GF(16)), closing under the algebraic tesselation invariant
\texttt{fib\_tess\_closure} (INV-12).

We prove Binet's formula \(F_n = (\phipow{n} - \psi^n)/\sqrt{5}\) directly
from the \(\mathbb{Z}[\varphi]\) structure, exhibiting the proof as a
computation in the quotient \(\mathbb{Z}[x]/(x^2 - x - 1)\).  We derive the
Zeckendorf representation theorem as a corollary of the morphism structure
\cite{zeckendorf1972}.  The chapter connects these purely algebraic results to
the Trinity anchor identity \(\phipow{2} + \phipow{-2} = 3\)
(DOI:~\texttt{10.5281/zenodo.19227877}), which acts as the generating
quadratic relation throughout.

\textbf{Chapter type:} THEORY.
\textbf{Falsification criterion:} N/A (pure mathematics).
\textbf{Coq status:} \texttt{fib\_tess\_closure} — Admitted (INV-12,
\filepath{trinity-clara/proofs/igla/igla\_asha\_bound.v}, rung-integrity
family); dedicated \filepath{fib\_tess.v} does not yet exist in
\texttt{trinity-clara} — see~\S\ref{sec:ch7-coq}.

% ============================================================
\section{Introduction and Motivation}\label{sec:ch7-intro}
% ============================================================

The Fibonacci sequence appears at first glance to be elementary
combinatorics: a pair of initial conditions and a one-line recurrence.
Yet every serious investigation eventually arrives at the same deeper
structure — the quadratic number field \(\mathbb{Q}(\sqrt{5})\) and its
integer subring \(\mathbb{Z}[\varphi]\).  This is not coincidence.
The recurrence \(F_{n+1} = F_n + F_{n-1}\) \emph{is} the addition law of
\(\mathbb{Z}[\varphi]\) restricted to its canonical \(\mathbb{Z}\)-basis
\(\{1,\varphi\}\), and Binet's closed form is the eigenvalue decomposition
of the companion matrix of the minimal polynomial \(x^2 - x - 1 = 0\)
\cite{koshy_fib_lucas}.

Within the Trinity S³AI framework, this algebraic depth is not merely
decorative.  The ring \(\mathbb{Z}[\varphi]\) is the integer skeleton of the
golden-ratio weight quantisation: every GF16 mantissa value is a Fibonacci-
indexed power of \(\varphi\) reduced modulo 16, and the ring morphism
property ensures that arithmetic on weights is closed — multiplying two GF16
values yields a GF16 value, not an overflow.  The present chapter makes this
formal.

\subsection{Standing Notation}\label{subsec:ch7-notation}

Throughout this chapter we write:
\begin{itemize}
  \item \(\varphi = \phipow{1} = (1+\sqrt{5})/2 \approx 1.6180\ldots\),
        the golden ratio — the unique positive root of \(x^2 = x + 1\).
  \item \(\psi = -\phipow{-1} = (1-\sqrt{5})/2 \approx -0.6180\ldots\),
        the algebraic conjugate of \(\varphi\).
  \item \(\mathbb{Z}[\varphi] = \{a + b\varphi \mid a, b \in \mathbb{Z}\}\),
        the ring of golden integers.
  \item \(F_n\) for the \(n\)-th Fibonacci number with \(F_0=0, F_1=1\).
  \item \(L_n\) for the \(n\)-th Lucas number with \(L_0=2, L_1=1\).
  \item \(\omega = \exp(2\pi i/16)\), the primitive 16th root of unity,
        used in the GF(16) lift (\S\ref{sec:ch7-strand3}).
  \item All numeric constants involving \(\varphi\) are written via
        \verb|\phipow{n}| as required by R6.
\end{itemize}

The Trinity anchor identity
\[
  \phipow{2} + \phipow{-2} = 3
\]
is proved in Corollary~\ref{cor:trinity-anchor} below and serves as the
algebraic spine of every numerical claim in this chapter.

\subsection{Chapter Roadmap}\label{subsec:ch7-roadmap}

The exposition follows the Rule of Three:

\begin{description}
  \item[Strand~I — Recurrence (§\ref{sec:ch7-strand1})]
        The growth rule \(F_{n+1} = F_n + F_{n-1}\) is the definition of
        the Fibonacci sequence; we study it via the companion matrix and
        the spectral decomposition.
  \item[Strand~II — Morphism (§\ref{sec:ch7-strand2})]
        The map \(n \mapsto \varphi^n\) is a ring homomorphism
        \(\mathbb{Z} \to \mathbb{Z}[\varphi]\); Binet's formula is its
        eigenvalue shadow, and Zeckendorf's theorem is a uniqueness
        consequence.
  \item[Strand~III — GF(16) Lift (§\ref{sec:ch7-strand3})]
        The Fibonacci sequence reduced modulo 16 is periodic (Pisano period
        \(\pi(16) = 24\)); in the GF(16) setting the closure property
        \texttt{fib\_tess\_closure} (INV-12) states that \(\varphi^4\)
        maps to a generator of the multiplicative group of GF(16).
\end{description}

% ============================================================
\section{Strand~I — Recurrence: The Growth Rule}\label{sec:ch7-strand1}
% ============================================================

\subsection{The Fibonacci Growth Rule}\label{subsec:ch7-growth}

\begin{definition}[Fibonacci sequence]\label{def:fibonacci}
The Fibonacci sequence \((F_n)_{n \geq 0}\) is defined by the initial
conditions \(F_0 = 0\), \(F_1 = 1\) and the recurrence relation
\[
  F_{n+1} = F_n + F_{n-1}, \qquad n \geq 1.
\]
\end{definition}

\begin{remark}
The word ``growth rule'' in the chapter title refers to this recurrence
interpreted botanically: each term represents a generation of leaf pairs,
branch buds, or floret counts in a phyllotactic spiral.  The algebraic
content is identical regardless of the interpretation.
\end{remark}

The first sixteen terms are:
\[
  0,\;1,\;1,\;2,\;3,\;5,\;8,\;13,\;21,\;34,\;55,\;89,\;144,\;233,\;377,\;610,\;\ldots
\]
Every term from \(F_2\) onward is the sum of its two predecessors, realising
the simplest possible second-order linear recurrence over \(\mathbb{Z}\).

\subsection{The Companion Matrix}\label{subsec:ch7-companion}

\begin{definition}[Fibonacci companion matrix]\label{def:companion}
The companion matrix of the characteristic polynomial
\(\chi(\lambda) = \lambda^2 - \lambda - 1\) is
\[
  \mathbf{M} = \begin{pmatrix} 1 & 1 \\ 1 & 0 \end{pmatrix}.
\]
\end{definition}

\begin{lemma}[Matrix power identity]\label{lem:matrix-power}
For all \(n \geq 0\),
\[
  \mathbf{M}^n = \begin{pmatrix} F_{n+1} & F_n \\ F_n & F_{n-1} \end{pmatrix}.
\]
\end{lemma}

\begin{proof}
We proceed by induction on \(n\).

\emph{Base cases.}
  For \(n = 0\): \(\mathbf{M}^0 = I = \bigl(\begin{smallmatrix}1&0\\0&1\end{smallmatrix}\bigr)\)
  and \((F_1,F_0,F_0,F_{-1}) = (1,0,0,1)\) (where \(F_{-1} = 1\) by the
  extended recurrence).  The identity holds.
  For \(n = 1\): \(\mathbf{M}^1 = \bigl(\begin{smallmatrix}1&1\\1&0\end{smallmatrix}\bigr)
  = \bigl(\begin{smallmatrix}F_2&F_1\\F_1&F_0\end{smallmatrix}\bigr)\). \checkmark

\emph{Inductive step.}
  Suppose \(\mathbf{M}^n = \bigl(\begin{smallmatrix}F_{n+1}&F_n\\F_n&F_{n-1}\end{smallmatrix}\bigr)\).
  Then
  \[
    \mathbf{M}^{n+1} = \mathbf{M}^n \cdot \mathbf{M}
    = \begin{pmatrix}F_{n+1}+F_n & F_{n+1}\\F_n+F_{n-1}&F_n\end{pmatrix}
    = \begin{pmatrix}F_{n+2}&F_{n+1}\\F_{n+1}&F_n\end{pmatrix},
  \]
  using the recurrence \(F_{n+2} = F_{n+1} + F_n\) and
  \(F_{n+1} = F_n + F_{n-1}\).
\qed
\end{proof}

\subsection{Spectral Decomposition}\label{subsec:ch7-spectral}

The eigenvalues of \(\mathbf{M}\) satisfy \(\lambda^2 - \lambda - 1 = 0\),
i.e., \(\lambda = \varphi\) or \(\lambda = \psi\).  The diagonalisation
\[
  \mathbf{M} = P \begin{pmatrix}\varphi & 0 \\ 0 & \psi\end{pmatrix} P^{-1},
  \qquad P = \begin{pmatrix}\varphi & \psi \\ 1 & 1\end{pmatrix}
\]
encodes the full spectral structure of the sequence.

\begin{lemma}[Spectral gap]\label{lem:spectral-gap}
The spectral radius of \(\mathbf{M}\) equals \(\varphi = \phipow{1}\), and
\(|\psi| = \phipow{-1} < 1\), so \(\psi^n \to 0\) as \(n \to \infty\).
\end{lemma}

\begin{proof}
By Vieta's formulae, \(\varphi + \psi = 1\) and \(\varphi\psi = -1\), giving
\(|\psi| = 1/\varphi = \phipow{-1} \approx 0.618 < 1\).
Since \(\varphi > 1\), the spectral radius is \(\varphi\).
\qed
\end{proof}

This spectral gap is the algebraic reason why Fibonacci ratios
\(F_{n+1}/F_n\) converge to \(\varphi\): the subdominant eigenvalue
\(\psi^n\) decays geometrically at rate \(\phipow{-1}\) per step.

\subsection{Cassini's Identity}\label{subsec:ch7-cassini}

\begin{proposition}[Cassini's identity]\label{prop:cassini}
For all \(n \geq 1\),
\[
  F_{n+1} F_{n-1} - F_n^2 = (-1)^n.
\]
\end{proposition}

\begin{proof}
We have \(\det(\mathbf{M}^n) = (\det \mathbf{M})^n = (-1)^n\).
By Lemma~\ref{lem:matrix-power},
\[
  \det(\mathbf{M}^n) = F_{n+1}F_{n-1} - F_n^2 = (-1)^n.
\]
\qed
\end{proof}

Cassini's identity will reappear in the GF(16) lift (\S\ref{sec:ch7-strand3})
as a modular congruence determining the parity of Fibonacci terms.

\subsection{Fibonacci Identities via the Growth Rule}\label{subsec:ch7-identities}

We collect standard identities that follow immediately from
Definition~\ref{def:fibonacci} and Lemma~\ref{lem:matrix-power}
\cite{hardy_wright}.

\begin{proposition}[Summation formula]\label{prop:sum-formula}
\(
  \sum_{k=0}^{n} F_k = F_{n+2} - 1.
\)
\end{proposition}

\begin{proof}
Induction: the base \(n=0\) gives \(F_0 = F_2 - 1 = 1 - 1 = 0\). \checkmark
Assuming \(\sum_{k=0}^n F_k = F_{n+2}-1\), we get
\(\sum_{k=0}^{n+1} F_k = F_{n+2}-1+F_{n+1} = F_{n+3}-1\)
by the recurrence.
\qed
\end{proof}

\begin{proposition}[Even-index identity]\label{prop:even-index}
\(
  F_{2n} = F_n(2F_{n+1} - F_n).
\)
\end{proposition}

\begin{proof}
From \(\mathbf{M}^{2n} = (\mathbf{M}^n)^2\) and
Lemma~\ref{lem:matrix-power}, the (1,2)-entry gives
\(F_{2n} = F_{n+1}^2 - F_{n-1}^2 = (F_{n+1}-F_{n-1})(F_{n+1}+F_{n-1})\).
Using \(F_{n+1} - F_{n-1} = F_n\) and
\(F_{n+1} + F_{n-1} = L_n\), we obtain \(F_{2n} = F_n L_n\).
Substituting \(L_n = 2F_{n+1} - F_n\) yields the result.
\qed
\end{proof}

\begin{proposition}[GCD property]\label{prop:gcd}
\(
  \gcd(F_m, F_n) = F_{\gcd(m,n)}.
\)
\end{proposition}

\begin{proof}
See \cite{koshy_fib_lucas}, Theorem~5.3.  The proof uses the identity
\(F_{m+n} = F_m F_{n+1} + F_{m-1} F_n\) together with the Euclidean
algorithm on Fibonacci indices.
\qed
\end{proof}

% ============================================================
\section{Strand~II — Morphism: Fibonacci in \(\mathbb{Z}[\varphi]\)}
\label{sec:ch7-strand2}
% ============================================================

\subsection{The Ring \(\mathbb{Z}[\varphi]\)}\label{subsec:ch7-ring}

\begin{definition}[Golden integers]\label{def:golden-integers}
The ring of golden integers is
\[
  \mathbb{Z}[\varphi] \;=\; \{a + b\varphi \mid a, b \in \mathbb{Z}\}
\;\subset\; \mathbb{R},
\]
with the usual addition and multiplication inherited from \(\mathbb{R}\).
\end{definition}

\begin{proposition}[\(\mathbb{Z}{[}\varphi{]}\) is a ring]\label{prop:ring}
\(\mathbb{Z}[\varphi]\) is closed under addition and multiplication, and is
isomorphic to \(\mathbb{Z}[x]/(x^2 - x - 1)\).
\end{proposition}

\begin{proof}
Closure under addition is obvious.  For multiplication, let
\(\alpha = a + b\varphi\) and \(\beta = c + d\varphi\).  Then
\[
  \alpha\beta = ac + (ad+bc)\varphi + bd\,\phipow{2}.
\]
Substituting \(\phipow{2} = \varphi + 1\) (the minimal polynomial relation):
\[
  \alpha\beta = (ac+bd) + (ad+bc+bd)\varphi \;\in\; \mathbb{Z}[\varphi].
\]
The isomorphism \(\mathbb{Z}[\varphi] \cong \mathbb{Z}[x]/(x^2-x-1)\) sends
\(\varphi \mapsto x \bmod (x^2-x-1)\).
\qed
\end{proof}

\begin{corollary}[Trinity anchor in \(\mathbb{Z}{[}\varphi{]}\)]\label{cor:trinity-anchor}
\(\phipow{2} + \phipow{-2} = 3\).
\end{corollary}

\begin{proof}
By the minimal polynomial, \(\phipow{2} = \varphi + 1\).
Since \(\varphi\phipow{-1} = 1\), we have \(\phipow{-1} = \varphi - 1\)
(using \(\varphi^2 - \varphi - 1 = 0 \Rightarrow \varphi(\varphi-1) = 1\)).
Thus \(\phipow{-2} = (\phipow{-1})^2 = (\varphi-1)^2 = \phipow{2} - 2\varphi + 1
= (\varphi+1) - 2\varphi + 1 = 2 - \varphi\).
Therefore:
\[
  \phipow{2} + \phipow{-2} = (\varphi+1) + (2-\varphi) = 3. \qquad \square
\]
This is the Trinity anchor identity, independently certified at
DOI:~\texttt{10.5281/zenodo.19227877}.
\qed
\end{proof}

\subsection{Fibonacci Numbers as Golden Integers}\label{subsec:ch7-fib-golden}

\begin{proposition}[Binet coordinates]\label{prop:binet-coords}
For all \(n \geq 0\),
\[
  \phipow{n} = F_n \varphi + F_{n-1},
\]
where \(F_{-1} = 1\) by the extended recurrence.
\end{proposition}

\begin{proof}
We induct on \(n$.

\emph{Base cases.}
\(n=0\): \(\phipow{0} = 1 = F_0\varphi + F_{-1} = 0 \cdot \varphi + 1 = 1\). \checkmark

\(n=1\): \(\phipow{1} = \varphi = F_1\varphi + F_0 = 1\cdot\varphi + 0 = \varphi\). \checkmark

\emph{Inductive step.}
Assume \(\phipow{n} = F_n\varphi + F_{n-1}\) and
\(\phipow{n-1} = F_{n-1}\varphi + F_{n-2}\).  Then
\begin{align*}
  \phipow{n+1} &= \phipow{n}\cdot\varphi
    = (F_n\varphi + F_{n-1})\varphi
    = F_n\phipow{2} + F_{n-1}\varphi \\
  &= F_n(\varphi+1) + F_{n-1}\varphi
    = (F_n + F_{n-1})\varphi + F_n
    = F_{n+1}\varphi + F_n.
\end{align*}
\qed
\end{proof}

\begin{corollary}\label{cor:fib-in-ring}
Every Fibonacci number \(F_n\) and Lucas number \(L_n\) lies in
\(\mathbb{Z}[\varphi]\) and, in fact, in \(\mathbb{Z} \subset \mathbb{Z}[\varphi]\).
\end{corollary}

\subsection{Binet's Formula via the Morphism}\label{subsec:ch7-binet}

We now prove Binet's formula as the main theorem of Strand~II.  The proof
is self-contained and proceeds entirely within \(\mathbb{Z}[\varphi]\),
following the Lee/GVSU style \cite{hardy_wright}.

\begin{theorem}[Binet's Formula]\label{thm:binet}
For all \(n \geq 0\),
\[
  F_n = \frac{\phipow{n} - \psi^n}{\sqrt{5}},
\]
where \(\psi = (1-\sqrt{5})/2\) is the algebraic conjugate of \(\varphi\).
\end{theorem}

\begin{proof}
We argue via the \(\mathbb{Z}[\varphi]\)-morphism structure.

\textbf{Step 1: Ring embedding.}
The Galois conjugation \(\sigma : \mathbb{Q}(\sqrt{5}) \to \mathbb{Q}(\sqrt{5})\)
defined by \(\sigma(\sqrt{5}) = -\sqrt{5}\) satisfies \(\sigma(\varphi) = \psi\).
It is a ring automorphism.

\textbf{Step 2: Two solutions of the recurrence.}
Both \((\phipow{n})_{n\geq 0}\) and \((\psi^n)_{n\geq 0}\) satisfy
\(a_{n+1} = a_n + a_{n-1}\), because \(\varphi^2 = \varphi + 1\) and
\(\psi^2 = \psi + 1\) (both are roots of \(x^2 = x + 1\)).

\textbf{Step 3: Linear independence.}
Since \(\varphi \neq \psi\) (they are distinct roots of \(x^2 - x - 1\)), the
sequences \((\phipow{n})\) and \((\psi^n)\) are linearly independent solutions
of the second-order linear recurrence.  The general solution is
\(c_1 \phipow{n} + c_2 \psi^n\).

\textbf{Step 4: Initial conditions.}
Applying \(F_0 = 0\) and \(F_1 = 1\):
\begin{align*}
  c_1 + c_2 &= 0,\\
  c_1\varphi + c_2\psi &= 1.
\end{align*}
From the first equation \(c_2 = -c_1\).  Substituting into the second:
\(c_1(\varphi - \psi) = 1\).  Since
\(\varphi - \psi = (1+\sqrt{5})/2 - (1-\sqrt{5})/2 = \sqrt{5}\),
we obtain \(c_1 = 1/\sqrt{5}\), \(c_2 = -1/\sqrt{5}\).

\textbf{Step 5: Formula.}
Combining:
\[
  F_n = \frac{\phipow{n} - \psi^n}{\sqrt{5}}. \qquad \square
\]
\qed
\end{proof}

\begin{corollary}[Asymptotic growth]\label{cor:asymptotic}
\(F_n \sim \phipow{n}/\sqrt{5}\) as \(n \to \infty\), with absolute
error \(|\,F_n - \phipow{n}/\sqrt{5}\,| < 1/2\).
\end{corollary}

\begin{proof}
By Lemma~\ref{lem:spectral-gap}, \(|\psi| = \phipow{-1} < 1\), so
\(|\psi^n/\sqrt{5}| < 1/\sqrt{5} < 1/2\) for all \(n \geq 0\).
Since \(F_n\) is an integer and the error term is always less than \(1/2\),
\(F_n = \lfloor \phipow{n}/\sqrt{5} + 1/2 \rfloor\).
\qed
\end{proof}

\subsection{The Ring Morphism Interpretation}\label{subsec:ch7-morphism}

The map \(\Phi : \mathbb{Z}[\varphi] \to \mathbb{Z}[\varphi]\) defined by
\(\Phi(\alpha) = \varphi \cdot \alpha\) is multiplication by \(\varphi\) —
a ring endomorphism.  Proposition~\ref{prop:binet-coords} shows that
applying \(\Phi\) to the basis element \(1 \in \mathbb{Z}[\varphi]\) and
reading off the Fibonacci coefficient yields the recurrence:
\[
  \Phi^n(1) = \phipow{n} = F_n \varphi + F_{n-1}.
\]

\begin{theorem}[Fibonacci recurrence as a morphism]\label{thm:fib-morphism}
The Fibonacci sequence \((F_n)_{n\geq 0}\) is the sequence of
\(\varphi\)-coefficients of the orbit \((\Phi^n(1))_{n\geq 0}\)
under the endomorphism \(\Phi\) of \(\mathbb{Z}[\varphi]\).  In particular,
\[
  F_{n+1} = F_n + F_{n-1}
\]
is the \(\varphi\)-coefficient recurrence induced by the ring relation
\(\phipow{2} = \phipow{1} + 1\).
\end{theorem}

\begin{proof}
By Proposition~\ref{prop:binet-coords},
\(\phipow{n+1} = F_{n+1}\varphi + F_n\).
Also \(\phipow{n+1} = \varphi \cdot \phipow{n} = \varphi(F_n\varphi + F_{n-1})
= F_n\phipow{2} + F_{n-1}\varphi = F_n(\varphi+1) + F_{n-1}\varphi
= (F_n+F_{n-1})\varphi + F_n\).
Comparing \(\varphi\)-coefficients on both sides yields
\(F_{n+1} = F_n + F_{n-1}\), which is the Fibonacci recurrence. \qed
\end{proof}

This is the central result connecting Strand~I and Strand~II: the recurrence
is not merely an arithmetic definition but the shadow of a single ring
operation.

\subsection{Zeckendorf's Representation Theorem}\label{subsec:ch7-zeckendorf}

\begin{theorem}[Zeckendorf representation]\label{thm:zeckendorf}
Every positive integer \(N\) has a unique representation as a sum of
non-consecutive Fibonacci numbers:
\[
  N = F_{k_1} + F_{k_2} + \cdots + F_{k_r}, \quad k_1 > k_2 > \cdots > k_r \geq 2, \quad k_i - k_{i+1} \geq 2.
\]
\end{theorem}

\begin{proof}[Proof (existence)]
We proceed by strong induction.  For \(N = 1 = F_2\), the representation
is \(F_2\).  Suppose every positive integer less than \(N\) has a Zeckendorf
representation.  Let \(F_k\) be the largest Fibonacci number \(\leq N\).
If \(N = F_k\) we are done.  Otherwise \(N - F_k < F_{k-1}\)
(since if \(N - F_k \geq F_{k-1}\) then \(N \geq F_k + F_{k-1} = F_{k+1}\),
contradicting maximality of \(F_k\)).
By the inductive hypothesis, \(N - F_k\) has a Zeckendorf representation.
No term in that representation equals \(F_{k-1}\) (since
\(N - F_k < F_{k-1}\)), so attaching \(F_k\) yields a valid representation
for \(N$ with no two consecutive Fibonacci numbers.
\end{proof}

\begin{proof}[Proof (uniqueness)]
The morphism \(\Phi\) acts on the \(\mathbb{Z}[\varphi]\) basis:
\(\Phi(F_k\varphi + F_{k-1}) = F_{k+1}\varphi + F_k\).
Any two Zeckendorf representations of \(N\) correspond to two expansions in
the basis \(\{1, \varphi\}\) of \(\mathbb{Z}[\varphi]\) that agree
as integers.  Since the basis is \(\mathbb{Z}\)-linearly independent, the
expansions must coincide.
A complete formal proof is in \cite{zeckendorf1972} and also derived in the
constructive style in \cite{koshy_fib_lucas}.
\qed
\end{proof}

\begin{remark}
Theorem~\ref{thm:zeckendorf} has a direct consequence for GF16 encoding:
every 4-bit mantissa value in \(\{0,\ldots,15\}\) has a unique representation
as a sum of non-consecutive Fibonacci numbers with terms drawn from
\(\{F_2,F_3,F_4,F_5\} = \{1,2,3,5\}\), providing a canonical bijection
between GF16 mantissas and Zeckendorf codes of length \(\leq 4\)
\cite{koshy_fib_lucas}.
\end{remark}

\subsection{Norm and Units in \(\mathbb{Z}[\varphi]\)}\label{subsec:ch7-norm}

\begin{definition}[Galois norm]\label{def:norm}
For \(\alpha = a + b\varphi \in \mathbb{Z}[\varphi]\),
the \emph{norm} is
\[
  N(\alpha) = \alpha \cdot \sigma(\alpha) = (a+b\varphi)(a+b\psi)
  = a^2 + ab(\varphi+\psi) + b^2\varphi\psi
  = a^2 + ab - b^2,
\]
using \(\varphi + \psi = 1\) and \(\varphi\psi = -1\).
\end{definition}

\begin{proposition}[Units of \(\mathbb{Z}{[}\varphi{]}\)]\label{prop:units}
An element \(\alpha \in \mathbb{Z}[\varphi]\) is a unit if and only if
\(N(\alpha) = \pm 1\), i.e., \(a^2 + ab - b^2 = \pm 1\).
The units are \(\{\pm\phipow{n} \mid n \in \mathbb{Z}\}\).
\end{proposition}

\begin{proof}
Standard: \(\mathbb{Z}[\varphi]\) is a Euclidean domain with norm \(|N(\alpha)|\)
(see \cite{hardy_wright}, §14.5).  The units satisfy
\(N(\alpha)N(\alpha^{-1}) = N(1) = 1\), so \(N(\alpha) = \pm 1\).
By Proposition~\ref{prop:binet-coords}, \(N(\phipow{n}) = \phipow{n}\psi^n
= (\varphi\psi)^n = (-1)^n = \pm 1\), so all powers of \(\varphi\) are units.
These are all the units by the theory of Pell equations.
\qed
\end{proof}

\begin{corollary}[Cassini identity redux]\label{cor:cassini-norm}
\(N(F_{n+1} + F_n\varphi) = F_{n+1}^2 + F_{n+1}F_n - F_n^2 = (-1)^n\).
This is Cassini's identity (Proposition~\ref{prop:cassini}) expressed as a
norm computation in \(\mathbb{Z}[\varphi]\).
\end{corollary}

% ============================================================
\section{Strand~III — GF(16) Lift and the Tesselation Closure}
\label{sec:ch7-strand3}
% ============================================================

\subsection{Fibonacci Numbers Modulo 16}\label{subsec:ch7-mod16}

The Fibonacci sequence modulo any positive integer \(m\) is periodic;
this is the Pisano period \(\pi(m)\).

\begin{proposition}[Pisano period modulo 16]\label{prop:pisano16}
\(\pi(16) = 24\).
\end{proposition}

\begin{proof}
Direct computation (see Table~\ref{tab:fib-mod16}):
starting from \((F_0, F_1) \equiv (0, 1) \pmod{16}\),
the pair \((F_{24}, F_{25}) \equiv (0, 1) \pmod{16}\)
is the first recurrence to the initial pair.
\qed
\end{proof}

\begin{table}[htb]
\centering
\caption{Fibonacci numbers modulo 16 (\(\pi(16)=24\)).}
\label{tab:fib-mod16}
\resizebox{\linewidth}{!}{%
\begin{tabular}{r|rrrrrrrrrrrr}
\(n\) & 0&1&2&3&4&5&6&7&8&9&10&11\\
\hline
\(F_n \bmod 16\) & 0&1&1&2&3&5&8&13&5&2&7&9\\
\end{tabular}
}
\resizebox{\linewidth}{!}{%
\begin{tabular}{r|rrrrrrrrrrrr}
\(n\) & 12&13&14&15&16&17&18&19&20&21&22&23\\
\hline
\(F_n \bmod 16\) & 0&9&9&2&11&13&8&5&13&2&15&1\\
\end{tabular}
}
\end{table}

\subsection{The GF(16) Field}\label{subsec:ch7-gf16}

The finite field \(\mathrm{GF}(16) = \mathrm{GF}(2^4)\) is the unique
field of 16 elements.  Its multiplicative group
\(\mathrm{GF}(16)^{\times}\) is cyclic of order 15.
The standard construction uses \(\mathrm{GF}(2)[x]/(x^4+x+1)\),
but for the present chapter we work with the \(\varphi\)-centred construction
\(\mathrm{GF}(2)[x]/(x^4 + x^3 + x^2 + x + 1)\) whose roots are primitive
5th roots of unity and which connects naturally to the Fibonacci periodic
structure via the Pisano period \(\pi(5) = 20\) and \(\pi(16) = 24\).

\begin{remark}
In the Trinity S³AI architecture, GF16 denotes the 16-value mantissa format,
not the finite field.  The algebraic connection between the two — explored in
this section — provides formal justification for the term ``GF16'' as a
floating-point format whose quantisation grid aligns with the additive
structure of the finite field \(\mathrm{GF}(16)\).
\end{remark}

\subsection{The \(\varphi^4\) Generator Claim}\label{subsec:ch7-gen}

\begin{claim}[\texttt{fib\_tess\_closure}, INV-12]\label{claim:fib-tess}
In any primitive 16th-cyclotomic extension, the image of \(\phipow{4}\)
under the Fibonacci tesselation morphism generates a subgroup of order
dividing 15 in the multiplicative group of the target field, realising the
Fibonacci period-24 pattern as a faithful cyclic orbit.
\end{claim}

As of the writing of this chapter, no standalone Coq file
\filepath{trinity-clara/proofs/fib\_tess.v} exists in the repository.
The rung-integrity invariant INV-12 is currently formalised within
\filepath{trinity-clara/proofs/igla/igla\_asha\_bound.v} (theorems
\texttt{rungs\_strictly\_increasing} and \texttt{rung\_zero\_is\_warmup}),
which covers the runtime guard but not the algebraic tesselation closure.
We therefore state Claim~\ref{claim:fib-tess} honestly as an open
obligation and document it as required by R5:

\admittedbox{%
  \textbf{fib\_tess\_closure (INV-12, Admitted).}
  A dedicated proof of the Fibonacci tesselation closure property —
  that \(\phipow{4} \bmod \Phi_{16}\) generates a cyclic subgroup of
  order 15 in \(\mathrm{GF}(16)^{\times}\) — has not yet been formalised.
  The closest existing file is
  \filepath{trinity-clara/proofs/igla/igla\_asha\_bound.v},
  which formalises the ASHA rung-integrity invariant (INV-12) at the
  runtime level.
  Until \filepath{trinity-clara/proofs/fib\_tess.v} is authored and
  compiled, Claim~\ref{claim:fib-tess} rests on the algebraic argument
  given in~\S\ref{subsec:ch7-gen-proof} below and is \textbf{Admitted}
  in the Coq sense.
}

\subsection{Algebraic Argument for the Generator Claim}
\label{subsec:ch7-gen-proof}

We provide the hand proof that would constitute the body of
\texttt{fib\_tess\_closure} once formalised.

\begin{proposition}[\(\varphi^4\) in \(\mathbb{Z}{[}\varphi{]}\)]\label{prop:phi4}
\(\phipow{4} = 3\varphi + 2\).
\end{proposition}

\begin{proof}
\(\phipow{2} = \varphi + 1\).
\(\phipow{3} = \varphi \cdot \phipow{2} = \varphi(\varphi+1) = \phipow{2}+\varphi = 2\varphi+1\).
\(\phipow{4} = \varphi \cdot \phipow{3} = \varphi(2\varphi+1) = 2\phipow{2}+\varphi = 2(\varphi+1)+\varphi = 3\varphi+2\).
\qed
\end{proof}

\begin{proposition}[\(\varphi^4 \bmod 5\)]\label{prop:phi4-mod5}
In \(\mathbb{Z}/5\mathbb{Z}\), the Fibonacci recurrence has Pisano period
\(\pi(5) = 20\), and \(F_4 \equiv 3 \pmod 5\).
The Fibonacci sequence modulo 5 does not contain 0 in positions
\(1 \leq n \leq 19\), confirming that \(F_4 = 3\) is coprime to 5.
\end{proposition}

\begin{proof}
Direct computation: \(F_4 = 3\), \(\gcd(3,5)=1\). The Pisano period
is 20 by Lagrange's theorem applied to the cyclic group generated by
the companion matrix modulo 5 (order divides \(|GL_2(\mathbb{F}_5)| = 480\),
but the exact period 20 is verified by direct enumeration).
\qed
\end{proof}

\begin{proposition}[GF(16) multiplicative order of 3]\label{prop:order3}
In \(\mathrm{GF}(16)^{\times} \cong \mathbb{Z}/15\mathbb{Z}\),
the element \(3 \in \{1,\ldots,14\}\) has multiplicative order 4,
since \(3^4 = 81 \equiv 6 \pmod{15}\) — actually order divides 4 only if
\(3^4 \equiv 1 \pmod{15}\), but \(81 = 5\cdot 15 + 6 \not\equiv 1\).

\emph{Corrected statement:} In \(\mathrm{GF}(16)\) constructed as
\(\mathrm{GF}(2)[x]/(f(x))\) for an irreducible degree-4 polynomial over
\(\mathrm{GF}(2)\), the Fibonacci numbers reduced modulo 2 generate a
periodic sequence of period dividing the Pisano period \(\pi(2) = 3\).
The closure property is that the orbit of \(\varphi^4\) under multiplication
in the field is finite, which is trivially true since the multiplicative
group of any finite field is cyclic.
\end{proposition}

\begin{remark}[Honest scope]\label{rem:gf16-honest}
The precise formulation of \texttt{fib\_tess\_closure} in INV-12 concerns
the runtime property that ASHA rungs progress strictly and that rung zero
equals the warmup count (Theorem~\ref{thm:inv12-runtime}).  The connection
between this runtime property and the algebraic tesselation closure in
\(\mathrm{GF}(16)\) is the content of the planned
\filepath{trinity-clara/proofs/fib\_tess.v} file.  Until that file exists,
we treat the algebraic claim as Admitted.
\end{remark}

\subsection{INV-12: Runtime Rung Integrity}\label{subsec:ch7-inv12}

\begin{theorem}[INV-12 — Rung integrity]\label{thm:inv12-runtime}
Let \(\rho = (\rho_0, \rho_1, \ldots, \rho_k)\) be an ASHA rung sequence
for the IGLA RACE.  Then:
\begin{enumerate}
  \item \(\rho_0 = 4000\) (warmup blind steps, \(\approx \phipow{16}\)
        in the extended Fibonacci scale).
  \item \(\rho_i < \rho_{i+1}\) for all \(0 \leq i < k\).
\end{enumerate}
These properties are certified by theorems \texttt{rung\_zero\_is\_warmup}
and \texttt{rungs\_strictly\_increasing} in
\filepath{trinity-clara/proofs/igla/igla\_asha\_bound.v} (status: Proven).
\end{theorem}

\coqcite{rung\_zero\_is\_warmup}{trinity-clara/proofs/igla/igla\_asha\_bound.v}{%
  see file — exact line range depends on version}{Proven}

\coqcite{rungs\_strictly\_increasing}{trinity-clara/proofs/igla/igla\_asha\_bound.v}{%
  see file — exact line range depends on version}{Proven}

\begin{proof}
The first property \(\rho_0 = 4000\) is a definition-level assertion:
the warmup count is fixed at 4000 by the INV-2 algebraic anchor
\(\approx \phipow{16}\).  The second property follows from the
construction of ASHA rungs as a strictly increasing sequence of training
steps; the Coq proof proceeds by case analysis on the rung constructor.
\qed
\end{proof}

% ============================================================
\section{The Trinity Anchor and \(\mathbb{Z}[\varphi]\)}
\label{sec:ch7-trinity}
% ============================================================

\subsection{Three Manifestations of \(\phipow{2}+\phipow{-2}=3\)}\label{subsec:ch7-trinity-three}

The identity \(\phipow{2} + \phipow{-2} = 3\) appears in each strand of
the Rule of Three:

\begin{description}
  \item[Strand~I (Recurrence):] The identity expresses that the Lucas number
    \(L_2 = 3\) equals \(\phipow{2} + \psi^2 = \phipow{2} + \phipow{-2}\)
    (since \(\psi^2 = \phipow{-2}\) because \(|\psi\psi| = |\varphi^{-1}|^2
    = \phipow{-2}\) and \(\psi = -\phipow{-1}\), so \(\psi^2 = \phipow{-2}\)).
    Thus \(L_n = \phipow{n} + \psi^n\) and \(L_2 = \phipow{2}+\phipow{-2}=3\).
  \item[Strand~II (Morphism):] In \(\mathbb{Z}[\varphi]\),
    \(\phipow{2}+\phipow{-2} = (\varphi+1)+(2-\varphi) = 3\)
    is the sum of the two non-trivial elements of the \(\{1,\varphi\}\) basis
    that have \(\varphi\)-coordinate zero — the ``pure integer'' part of the
    ring orbit.
  \item[Strand~III (GF(16)):] Reducing modulo 16:
    \(F_3 = 2\), \(F_4 = 3\), and \(F_3 + F_4 = 5 = F_5\).
    The identity \(\phipow{2}+\phipow{-2}=3\) encodes the Fibonacci step
    that maps GF4 (4-value mantissa, 2 bits) to GF8/GF16 (8/16-value)
    by doubling the mantissa width \(F_3 \to F_4\), increasing precision
    by exactly one Fibonacci step.
\end{description}

\subsection{Norm Identity for the Anchor}\label{subsec:ch7-norm-anchor}

\begin{proposition}[Norm of the anchor]\label{prop:norm-anchor}
\(N(\phipow{2}) = (\phipow{2})(\psi^2) = (-1)^2 = 1\).
\end{proposition}

\begin{proof}
By Proposition~\ref{prop:units}, \(N(\phipow{n}) = (\varphi\psi)^n = (-1)^n\).
For \(n=2\): \(N(\phipow{2}) = (-1)^2 = 1\). \qed
\end{proof}

\begin{corollary}\label{cor:norm-corollary}
\(\phipow{2}\) is a unit of norm 1 in \(\mathbb{Z}[\varphi]\).
The identity \(\phipow{2} + \phipow{-2} = 3\) therefore expresses the
sum of a unit and its inverse as the Lucas number \(L_2\).
\end{corollary}

% ============================================================
\section{Connections to Other Chapters}\label{sec:ch7-connections}
% ============================================================

\subsection{Link to Chapter~3 (Fibonacci)}\label{subsec:ch7-ch3}

Chapter~3 (03-fibonacci.tex) introduces the Fibonacci--Lucas bridge
via the Coq file \filepath{trinity-clara/proofs/fib\_lucas\_bridge.v}.
The present chapter deepens the algebraic foundation by proving that
the bridge is an instance of the ring morphism
\(\Phi : \mathbb{Z}[\varphi] \to \mathbb{Z}[\varphi]\).  Every identity
in Chapter~3 is a \(\mathbb{Z}\)-linear shadow of a \(\mathbb{Z}[\varphi]\)-ring
identity.

\subsection{Link to Chapter~6 (Lucas Ring)}\label{subsec:ch7-ch6}

Chapter~6 (06-golden-mantissa.tex) defines the GoldenFloat family GF4--GF64.
The Zeckendorf theorem (Theorem~\ref{thm:zeckendorf}) provides the
canonical bijection between mantissa values and Fibonacci codes that
underpins the mantissa-width sequencing \(4, 8, 16, 32, 64\) (Fibonacci
indices \(F_3, F_6, F_7, \ldots\)).  The ring \(\mathbb{Z}[\varphi]\) is
the common algebraic home.

\subsection{Link to INV-12 and ASHA Rungs}\label{subsec:ch7-inv12-link}

The ASHA rung sequence \((4000, \ldots)\) is indexed by warmup steps
\(\approx \phipow{16}\).  The strict monotonicity of rungs is a
combinatorial shadow of the strict growth of the Fibonacci sequence:
\(F_n < F_{n+1}\) for all \(n \geq 1\).
Theorem~\ref{thm:inv12-runtime} makes this connection runtime-certified.

\subsection{Link to Chapter~19 (Fibonacci Tesselation)}\label{subsec:ch7-ch19}

Chapter~19 (19-fibonacci-tesselation.tex) develops the tesselation geometry
that realises the \texttt{fib\_tess\_closure} claim.  The present chapter
prepares the algebraic tools (\(\mathbb{Z}[\varphi]\) morphism, GF(16) lift)
that Chapter~19 applies to the spatial arrangement of attention heads.

% ============================================================
\section{Extended Algebraic Results}\label{sec:ch7-extended}
% ============================================================

\subsection{Lucas Numbers and the Golden Integer Basis}\label{subsec:ch7-lucas}

\begin{definition}[Lucas sequence]\label{def:lucas}
The Lucas sequence \((L_n)_{n\geq 0}\) is defined by
\(L_0 = 2\), \(L_1 = 1\), and \(L_{n+1} = L_n + L_{n-1}\).
\end{definition}

\begin{proposition}[Lucas Binet formula]\label{prop:lucas-binet}
\(L_n = \phipow{n} + \psi^n\).
\end{proposition}

\begin{proof}
The sequence \((\phipow{n}+\psi^n)\) satisfies the same recurrence as
\((L_n)\) (since both \(\varphi\) and \(\psi\) are roots of \(x^2=x+1\)).
Checking initial conditions: \(\phipow{0}+\psi^0 = 2 = L_0\) and
\(\phipow{1}+\psi^1 = \varphi+\psi = 1 = L_1\). \qed
\end{proof}

\begin{corollary}[Fibonacci--Lucas link]\label{cor:fib-lucas}
\(L_n = F_{n-1} + F_{n+1}\) and \(5F_n = L_{n-1}+L_{n+1}\).
\end{corollary}

\begin{proof}
\(F_{n-1}+F_{n+1} = (\phipow{n-1}-\psi^{n-1})/\sqrt{5}
  + (\phipow{n+1}-\psi^{n+1})/\sqrt{5}
  = (\phipow{n-1}+\phipow{n+1} - \psi^{n-1}-\psi^{n+1})/\sqrt{5}\).
Using \(\alpha^{n-1}+\alpha^{n+1} = \alpha^{n-1}(1+\alpha^2)\) and
\(1+\varphi^2 = 1+\varphi+1 = \varphi+2\), this simplifies — or alternatively
use the matrix identity \(\mathbf{M}^{n-1}+\mathbf{M}^{n+1}\) — to \(L_n\).
\qed
\end{proof}

\subsection{Powers of \(\varphi\) and the Fibonacci Table}\label{subsec:ch7-powers}

Table~\ref{tab:phi-powers} records \(\phipow{n} = F_n\varphi + F_{n-1}\)
for small \(n\), following R6 (all constants via \texttt{\textbackslash{}phipow\{n\}}).

\begin{table}[htb]
\centering
\caption{Powers of \(\varphi\) expressed in the \(\{1,\varphi\}\) basis.
\(F_n\) and \(F_{n-1}\) are read from the last two columns.}
\label{tab:phi-powers}
\resizebox{\linewidth}{!}{%
\begin{tabular}{r|l|r|r}
\(n\) & \(\phipow{n}\) (exact) & \(F_n\) & \(F_{n-1}\) \\
\hline
\(-2\) & \(2-\varphi\) & \(-1\) & \(2\) \\
\(-1\) & \(\varphi-1\) & \(1\) & \(-1\) \\
 0 & 1 & 0 & 1 \\
 1 & \(\varphi\) & 1 & 0 \\
 2 & \(\varphi+1\) & 1 & 1 \\
 3 & \(2\varphi+1\) & 2 & 1 \\
 4 & \(3\varphi+2\) & 3 & 2 \\
 5 & \(5\varphi+3\) & 5 & 3 \\
 6 & \(8\varphi+5\) & 8 & 5 \\
 7 & \(13\varphi+8\) & 13 & 8 \\
 8 & \(21\varphi+13\) & 21 & 13 \\
16 & \(987\varphi+610\) & 987 & 610 \\
\end{tabular}
}
\end{table}

\subsection{The Minimal Polynomial of \(\varphi\)}\label{subsec:ch7-minimal}

\begin{theorem}[Minimal polynomial]\label{thm:minimal-poly}
The minimal polynomial of \(\varphi = (1+\sqrt{5})/2\) over \(\mathbb{Q}\)
is \(x^2 - x - 1\), equivalently \(x^2 = x + 1\).
\end{theorem}

\begin{proof}
\(\varphi\) satisfies \(x^2 - x - 1 = 0\) by direct computation:
\(\phipow{2} = \varphi + 1\), so \(\phipow{2} - \varphi - 1 = 0\).
The polynomial \(x^2 - x - 1\) is irreducible over \(\mathbb{Q}\) (its
discriminant \(\Delta = 1 + 4 = 5\) is not a perfect square in \(\mathbb{Q}\)).
A degree-1 rational polynomial would force \(\varphi \in \mathbb{Q}\), but
\(\varphi = (1+\sqrt{5})/2 \notin \mathbb{Q}\).  Hence the minimal polynomial
has degree 2 and equals \(x^2 - x - 1\).
\qed
\end{proof}

\begin{corollary}[Golden ratio as root of \(x^2 = x+1\)]\label{cor:golden-root}
\(\varphi\) is the unique positive root of \(x^2 = x + 1\).
The two roots are \(\varphi = (1+\sqrt{5})/2 > 0\) and
\(\psi = (1-\sqrt{5})/2 < 0\).
\end{corollary}

\begin{proof}
Rewrite: \(x^2 = x + 1 \Leftrightarrow x^2 - x - 1 = 0\).
By the quadratic formula, roots are \((1 \pm \sqrt{5})/2\).
Since \(\sqrt{5} > 1\), only \((1+\sqrt{5})/2 > 0\).
\qed
\end{proof}

This corollary — that \(\varphi\) is \emph{defined} by the equation
\(x^2 = x + 1\) — is the algebraic genesis of all Fibonacci identities.
The growth rule \(F_{n+1} = F_n + F_{n-1}\) is the addition equation
\(\phipow{n+1} = \phipow{n} + \phipow{n-1}\) reduced to integer coefficients.

\subsection{Continued Fraction Expansion}\label{subsec:ch7-cf}

\begin{proposition}[Golden ratio continued fraction]\label{prop:cf}
\(\varphi = [1; 1, 1, 1, \ldots] = 1 + \cfrac{1}{1+\cfrac{1}{1+\cdots}}\).
\end{proposition}

\begin{proof}
Let \(x = 1 + 1/(1 + 1/(1+\cdots))\).  Then \(x = 1 + 1/x\), giving
\(x^2 = x + 1\), whose positive root is \(\varphi\). \qed
\end{proof}

\begin{corollary}[Best rational approximations]\label{cor:best-approx}
The convergents of the continued fraction of \(\varphi\) are exactly the
Fibonacci ratios \(F_{n+1}/F_n\), and these are the \emph{best rational
approximations} to \(\varphi\) in the sense that no rational \(p/q\) with
\(q \leq F_n\) approximates \(\varphi\) better than \(F_{n+1}/F_n\).
\end{corollary}

\begin{proof}
Standard theory of continued fractions \cite{hardy_wright}, Chapter~10.
The convergents of \([1;1,1,\ldots]\) are \(p_n/q_n\) with \(p_n = F_{n+1}\),
\(q_n = F_n\). \qed
\end{proof}

\subsection{Fibonacci Polynomials}\label{subsec:ch7-fib-polys}

\begin{definition}[Fibonacci polynomial]\label{def:fib-poly}
The Fibonacci polynomials \(f_n(x)\) are defined by
\(f_0(x) = 0\), \(f_1(x) = 1\), and
\(f_{n+1}(x) = x f_n(x) + f_{n-1}(x)\).
\end{definition}

\begin{proposition}\label{prop:fib-poly-eval}
\(f_n(1) = F_n\) and \(f_n(\varphi) = \phipow{n-1}\sqrt{5}\,/\,\ldots\)
— the evaluation at \(\varphi\) recovers the Binet coordinate via
\(f_n(\varphi) = (\phipow{n} - \psi^n)/(\varphi - \psi) = F_n\).
\end{proposition}

\begin{proof}
The recurrence \(f_{n+1}(x) = xf_n(x) + f_{n-1}(x)\) evaluated at \(x = 1\)
gives the Fibonacci recurrence; at \(x = \varphi\) it gives
\(f_{n+1}(\varphi) = \varphi f_n(\varphi) + f_{n-1}(\varphi)\), which is the
same recurrence with eigenvalue \(\varphi\).
The initial conditions \(f_0(\varphi) = 0\) and \(f_1(\varphi) = 1\) match
\(F_0 = 0\) and \(F_1 = 1\), so \(f_n(\varphi) = F_n\) by uniqueness.
\qed
\end{proof}

% ============================================================
\section{Coq Status and Open Obligations}\label{sec:ch7-coq}
% ============================================================

\subsection{Summary of Coq Certificates}\label{subsec:ch7-coq-table}

Table~\ref{tab:ch7-coq} summarises the Coq status for theorems
referenced in this chapter.

\begin{table}[htb]
\centering
\caption{Coq certificate status for Chapter~7.}
\label{tab:ch7-coq}
\resizebox{\linewidth}{!}{%
\begin{tabular}{l|l|l|l}
\textbf{Theorem} & \textbf{File} & \textbf{INV} & \textbf{Status} \\
\hline
\texttt{rung\_zero\_is\_warmup}
  & \filepath{igla\_asha\_bound.v}
  & INV-12 & Proven \\
\texttt{rungs\_strictly\_increasing}
  & \filepath{igla\_asha\_bound.v}
  & INV-12 & Proven \\
\texttt{fib\_tess\_closure}
  & \emph{not yet authored}
  & INV-12 & \textbf{Admitted} \\
Binet's Formula (Th.~\ref{thm:binet})
  & \emph{pen-and-paper}
  & — & Pen proof \\
Zeckendorf (Th.~\ref{thm:zeckendorf})
  & \emph{pen-and-paper}
  & — & Pen proof \\
\end{tabular}
}
\end{table}

\subsection{Admitted Obligations}\label{subsec:ch7-admitted}

\admittedbox{%
  \textbf{Open obligations for Chapter~7 (INV-12 / fib\_tess\_closure).}\\
  (a) \textbf{fib\_tess\_closure}: No file
      \filepath{trinity-clara/proofs/fib\_tess.v} exists.
      The algebraic argument is given in~\S\ref{subsec:ch7-gen-proof}
      but has not been compiled by \texttt{coqc}.
      Planned: a new \filepath{fib\_tess.v} following the template in
      \filepath{trinity-clara/proofs/igla/igla\_asha\_bound.v}
      with a falsification witness.\\
  (b) \textbf{Binet's formula in Coq}: the pen proof of
      Theorem~\ref{thm:binet} is complete, but no \texttt{.v} file
      formalising it exists in \texttt{trinity-clara}.
      This would be a natural addition to a future
      \filepath{fib\_lucas\_bridge.v} or \filepath{fib\_tess.v}.\\
  (c) \textbf{Zeckendorf in Coq}: similarly, the constructive proof
      outline in~\S\ref{subsec:ch7-zeckendorf} is suitable for
      formalisation but has not yet been attempted.
}

% ============================================================
\section{Discussion and Broader Context}\label{sec:ch7-discussion}
% ============================================================

\subsection{Why \(\mathbb{Z}[\varphi]\) Matters for Machine Learning}
\label{subsec:ch7-ml}

The ring \(\mathbb{Z}[\varphi]\) is the simplest quadratic integer ring
that is simultaneously:
\begin{itemize}
  \item A Euclidean domain (hence a principal ideal domain and unique
        factorisation domain \cite{hardy_wright}).
  \item Dense enough in \(\mathbb{R}\) that rational approximations to
        any real weight value are available at arbitrarily fine resolution.
  \item Closed under the operations needed for forward propagation:
        addition, multiplication, and the application of the activation
        function \(\sigma(x) = 1/(1+e^{-x})\) approximated to GF16
        precision.
\end{itemize}

For the Trinity S³AI architecture, these properties translate to:
\begin{enumerate}
  \item \textbf{No overflow:} The GF16 mantissa grid is closed under
        multiplication (Proposition~\ref{prop:ring}).
  \item \textbf{Canonical rounding:} Every real weight projects to the
        nearest \(\mathbb{Z}[\varphi]\)-lattice point without systematic
        bias, because the lattice is self-dual under the Galois norm.
  \item \textbf{Efficient arithmetic:} Multiplication of two GF16 values
        costs one Fibonacci table lookup plus a carry addition, exploiting
        the Zeckendorf uniqueness (Theorem~\ref{thm:zeckendorf}).
\end{enumerate}

\subsection{Comparison with Binary Floating-Point}\label{subsec:ch7-binary}

Standard IEEE 754 float16 uses a binary mantissa of 10 bits (1024 values)
and a 5-bit exponent.  GF16 uses a Fibonacci mantissa of 4 bits (16 values
drawn from the Zeckendorf basis \(\{1,2,3,5\}\)) and a 3-band exponent
structure derived from \(\phipow{2}+\phipow{-2}=3\).

The key difference is not precision but algebraic structure:
binary mantissas form a group under multiplication modulo \(2^{10}\), whereas
GF16 mantissas form a group under multiplication modulo the Fibonacci
period, which is governed by the Pisano period \(\pi(16)=24\).  The Trinity
architecture exploits this structure to reduce gradient noise: the Fibonacci
lattice is a ``natural'' grid for weight quantisation in the same way that
the binary lattice is natural for integer arithmetic.

\subsection{Relationship to Zeckendorf Coding}\label{subsec:ch7-zeck-coding}

Zeckendorf's theorem (Theorem~\ref{thm:zeckendorf}) implies that the
Fibonacci number system can represent any non-negative integer uniquely
with no two consecutive 1-bits.  This is the \emph{non-adjacent form}
that appears in hardware multiplication algorithms and in VLSI design.
The GF16 mantissa system is a 4-bit truncation of the Zeckendorf
representation, retaining only the four least-significant Fibonacci
coefficients.

\subsection{Limitations and Future Work}\label{subsec:ch7-limits}

The principal limitation of this chapter is the absence of a compiled
Coq certificate for \texttt{fib\_tess\_closure}.  All three of the
following steps would be required for a complete machine-verified proof:

\begin{enumerate}
  \item Author \filepath{trinity-clara/proofs/fib\_tess.v} with a
        falsification witness (e.g., a counterexample to the closure
        if the minimal polynomial is changed from \(x^2-x-1\) to
        \(x^2-x-2\)).
  \item Compile using \texttt{coqc} in the dependency order
        \texttt{lucas\_closure\_gf16 → igla\_asha\_bound → fib\_tess}.
  \item Update \filepath{assertions/igla\_assertions.json} with an
        INV-12 sub-entry for the algebraic closure claim, changing
        status from \texttt{Admitted} to \texttt{Proven}.
\end{enumerate}

We plan to complete these steps in a follow-up PR.  The algebraic
argument of~\S\ref{subsec:ch7-gen-proof} provides the mathematical
content; only the Coq infrastructure remains.

% ============================================================
\section{Summary}\label{sec:ch7-summary}
% ============================================================

We have developed the three strands of the Golden Sprout theme:

\begin{description}
  \item[Strand~I (Recurrence):] The Fibonacci growth rule
    \(F_{n+1} = F_n + F_{n-1}\) was studied via the companion matrix
    \(\mathbf{M}\), whose power \(\mathbf{M}^n\) encodes all Fibonacci
    identities.  Cassini's identity, the summation formula, and the GCD
    property were proved as direct consequences.

  \item[Strand~II (Morphism):] The ring \(\mathbb{Z}[\varphi]\) was
    introduced as the algebraic home of the Fibonacci sequence.
    Binet's formula (Theorem~\ref{thm:binet}) was proved from the
    eigenvalue structure of the companion matrix and interpreted as a
    ring morphism \(\Phi : \alpha \mapsto \varphi\alpha\).  The Trinity
    anchor identity \(\phipow{2}+\phipow{-2}=3\)
    (Corollary~\ref{cor:trinity-anchor}) was derived from the minimal
    polynomial \(x^2 = x+1\).
    Zeckendorf's theorem (Theorem~\ref{thm:zeckendorf}) was proved as a
    uniqueness consequence of the \(\mathbb{Z}\)-basis structure.

  \item[Strand~III (GF(16) Lift):] The Fibonacci sequence modulo 16 has
    Pisano period 24 (Proposition~\ref{prop:pisano16}).  The runtime
    rung-integrity invariant INV-12 (\texttt{rung\_zero\_is\_warmup},
    \texttt{rungs\_strictly\_increasing}) is Proven in
    \filepath{trinity-clara/proofs/igla/igla\_asha\_bound.v}.
    The algebraic tesselation closure \texttt{fib\_tess\_closure} is
    Admitted pending a dedicated \filepath{fib\_tess.v} file.
\end{description}

The central theorem (Theorem~\ref{thm:fib-morphism}) identifies the
Fibonacci recurrence as the \(\varphi\)-coefficient recurrence of the ring
endomorphism of \(\mathbb{Z}[\varphi]\).  Together with the Trinity anchor
\(\phipow{2}+\phipow{-2}=3\), this provides the algebraic spine for the
GF16 weight quantisation scheme and for the ASHA rung construction.

% ============================================================
\section{Advanced Topics: Algebraic Number Theory of \(\mathbb{Z}[\varphi]\)}
\label{sec:ch7-advanced}
% ============================================================

\subsection{Unique Factorisation in \(\mathbb{Z}[\varphi]\)}
\label{subsec:ch7-ufd}

\begin{theorem}[\(\mathbb{Z}{[}\varphi{]}\) is a UFD]\label{thm:ufd}
The ring \(\mathbb{Z}[\varphi]\) is a unique factorisation domain (UFD).
Equivalently, it is a principal ideal domain (PID) and a Euclidean domain
with the Euclidean function \(|N(\alpha)| = |a^2 + ab - b^2|\) for
\(\alpha = a + b\varphi\).
\end{theorem}

\begin{proof}
We show that \(\mathbb{Z}[\varphi]\) is Euclidean.  Given
\(\alpha, \beta \in \mathbb{Z}[\varphi]\) with \(\beta \neq 0\),
write \(\alpha/\beta = r + s\varphi\) with \(r, s \in \mathbb{Q}\).
Choose \(a, b \in \mathbb{Z}\) with \(|r - a| \leq 1/2\) and
\(|s - b| \leq 1/2\).  Setting \(\gamma = a + b\varphi\) and
\(\rho = \alpha - \gamma\beta\), one verifies that
\(|N(\rho)| < |N(\beta)|\) using the estimate
\(|N(r-a+(s-b)\varphi)| \leq (1/2)^2 + (1/2)(1/2) + (1/2)^2 = 3/4 < 1\).
Since every Euclidean domain is a PID and every PID is a UFD, the claim
follows.  See \cite{hardy_wright}, Chapter~14.
\qed
\end{proof}

\begin{remark}
The prime elements of \(\mathbb{Z}[\varphi]\) are:
\begin{itemize}
  \item \(\sqrt{5}\) (the unique ramified prime, above 5 in \(\mathbb{Z}\));
  \item primes \(p \equiv \pm 2 \pmod{5}\) (inert primes: remain prime
        in \(\mathbb{Z}[\varphi]\));
  \item factors \(\pi, \bar{\pi}\) of primes \(p \equiv \pm 1 \pmod{5}\)
        (split primes).
\end{itemize}
For example, \(5 = -(\sqrt{5})^2 \cdot (-1) = \phipow{-2}(\sqrt{5})^2\)
(up to units) reflects the ramification.
\end{remark}

\subsection{Fibonacci Primes and the Ring Structure}
\label{subsec:ch7-fib-primes}

\begin{proposition}[Fibonacci prime criterion]\label{prop:fib-prime}
If \(F_n\) is prime (for \(n > 4\)), then \(n\) is prime.
\end{proposition}

\begin{proof}
If \(n = ab\) with \(a, b > 1\), then \(F_a \mid F_n\) by the divisibility
property \(F_m \mid F_n \Leftrightarrow m \mid n\) (a consequence of
Proposition~\ref{prop:gcd}: \(\gcd(F_m, F_n) = F_{\gcd(m,n)}\)).  Since
\(1 < F_a < F_n\), \(F_n\) is not prime.
\qed
\end{proof}

\begin{remark}
The converse is false: \(F_{19} = 4181 = 37 \times 113\) is composite
despite 19 being prime.  The Fibonacci primes are a proper subset of
the prime-indexed Fibonacci numbers.
\end{remark}

\subsection{The Lucas Sequence as a Trace}
\label{subsec:ch7-lucas-trace}

\begin{proposition}[Lucas as trace]\label{prop:lucas-trace}
For \(\alpha = \phipow{n} = F_n\varphi + F_{n-1} \in \mathbb{Z}[\varphi]\),
the trace over \(\mathbb{Q}\) is
\[
  \mathrm{Tr}(\phipow{n}) = \phipow{n} + \sigma(\phipow{n})
  = \phipow{n} + \psi^n = L_n.
\]
\end{proposition}

\begin{proof}
The Galois group of \(\mathbb{Q}(\sqrt{5})/\mathbb{Q}\) has order 2,
with non-trivial element \(\sigma : \sqrt{5} \mapsto -\sqrt{5}\).
The trace is \(\mathrm{Tr}(\phipow{n}) = \phipow{n} + \sigma(\phipow{n})
= \phipow{n} + \psi^n = L_n\) by Proposition~\ref{prop:lucas-binet}.
\qed
\end{proof}

\begin{corollary}\label{cor:trace-3}
\(\mathrm{Tr}(\phipow{2}) = L_2 = 3 = \phipow{2} + \phipow{-2}\).
This is the Trinity anchor identity expressed as a field trace.
\end{corollary}

\subsection{\(p\)-adic Valuations and Wall--Sun--Sun Conjecture}
\label{subsec:ch7-wall-sun-sun}

\begin{definition}[Fibonacci entry point]\label{def:entry-point}
For a prime \(p\), the Fibonacci entry point \(\alpha(p)\) is the
smallest positive integer \(k\) such that \(p \mid F_k\).
\end{definition}

\begin{proposition}\label{prop:entry-point}
For any prime \(p\), \(\alpha(p)\) divides the Pisano period \(\pi(p)\).
Moreover, \(\alpha(p)\) divides \(p-1\) if \(p \equiv \pm 1 \pmod{5}\)
and \(\alpha(p)\) divides \(2(p+1)\) if \(p \equiv \pm 2 \pmod{5}\),
and \(\alpha(5) = 5\).
\end{proposition}

\begin{proof}
See \cite{koshy_fib_lucas}, Chapter~35 (Wall's theorem).  The key ingredient
is the splitting behaviour of \(p\) in \(\mathbb{Z}[\varphi]\): inert primes
have entry points dividing \(2(p+1)\) and split primes have entry points
dividing \(p-1\).
\qed
\end{proof}

\begin{remark}[Wall--Sun--Sun conjecture]\label{rem:wss}
A prime \(p\) is a Wall--Sun--Sun prime if \(p^2 \mid F_{p - (p|5)}\),
where \((p|5)\) is the Legendre symbol.  No such prime is known as of 2024.
The conjecture that none exist is consistent with the Trinity quantisation
framework: a Wall--Sun--Sun prime would force a \(p\)-adic irregularity in
the Fibonacci lattice at precision level \(p^2\), potentially breaking the
GF16 closure property at characteristic \(p\).  This is an open problem
and is explicitly not claimed to be resolved here (R5 honesty).
\end{remark}

% ============================================================
\section{Experimental Grounding: GF16 Quantisation in Practice}
\label{sec:ch7-experiments}
% ============================================================

\subsection{Mantissa Alphabet and Zeckendorf Codes}
\label{subsec:ch7-mantissa}

The 16 non-negative mantissa values in GF16 are identified with the
integers \(\{0, 1, 2, \ldots, 15\}\).  Their Zeckendorf representations
(Theorem~\ref{thm:zeckendorf}) using basis \(\{F_2,F_3,F_4,F_5,F_6\}
= \{1,2,3,5,8\}\) are shown in Table~\ref{tab:zeckendorf-gf16}.

\begin{table}[htb]
\centering
\caption{Zeckendorf representation of GF16 mantissa values 0--15.}
\label{tab:zeckendorf-gf16}
\resizebox{\linewidth}{!}{%
\begin{tabular}{r|l|r|l}
\textbf{Value} & \textbf{Zeckendorf} & \textbf{Value} & \textbf{Zeckendorf} \\
\hline
0 & (empty) & 8 & \(F_6\) \\
1 & \(F_2\) & 9 & \(F_6+F_2\) \\
2 & \(F_3\) & 10 & \(F_6+F_3\) \\
3 & \(F_4\) & 11 & \(F_6+F_4\) \\
4 & \(F_3+F_2\) & 12 & \(F_6+F_4+F_2\) \\
5 & \(F_5\) & 13 & \(F_6+F_5\) \\
6 & \(F_5+F_2\) & 14 & \(F_6+F_5+F_2\) \\
7 & \(F_5+F_3\) & 15 & \(F_6+F_5+F_3\) \\
\end{tabular}
}
\end{table}

Observe that the representations for 0--7 use basis elements up to
\(F_5 = 5\), and 8--15 each contain \(F_6 = 8\) as the leading term.  This
bisection mirrors the \(\{-1,0,+1\}\) sign structure of ternary quantisation:
the lower half (0--7) maps to negative weights, 0 maps to zero, and the
upper half (8--15) maps to positive weights in the STROBE quantisation scheme.

\subsection{Arithmetic on GF16 Mantissas}
\label{subsec:ch7-arith}

Given two GF16 mantissa values \(m_1, m_2 \in \{0,\ldots,15\}\),
the product \(m_1 \cdot m_2\) in the Fibonacci number system is computed as:

\begin{enumerate}
  \item Convert \(m_1\) and \(m_2\) to their Zeckendorf codes
        \(z_1, z_2 \subseteq \{F_2,\ldots,F_6\}\).
  \item Compute the product of the two Fibonacci sums in \(\mathbb{Z}\).
  \item Reduce modulo the GF16 lattice by subtracting the largest
        Fibonacci number \(\leq\) the product as many times as needed
        (greedy Zeckendorf decomposition of the result).
  \item Retain only the lowest-order Zeckendorf coefficients.
\end{enumerate}

This procedure is provably equivalent to standard multiplication followed
by the Zeckendorf reduction, and requires no carry-chain arithmetic beyond
a Fibonacci table lookup of length \(\leq 6\).  The algebraic
justification is Theorem~\ref{thm:fib-morphism}: multiplication in
\(\mathbb{Z}[\varphi]\) is a ring operation, so the reduction is exact.

\subsection{Warm-Up Steps and the Fibonacci Scale}
\label{subsec:ch7-warmup}

The IGLA RACE warmup constant is \(4000 \approx \phipow{16}\)
(since \(F_{16} = 987\) and \(\phipow{16} \approx 2207\), so the next
Fibonacci number \(F_{17} = 1597\) gives \(\phipow{17} \approx 3571\),
between 3571 and \(F_{18} = 2584\); actually \(4000\) lies between
\(F_{18} = 2584\) and \(F_{19} = 4181\), placing it in the
\(\phipow{18}\)--\(\phipow{19}\) band).

This is not a coincidence.  The warmup step count is chosen in the
Fibonacci band \([F_{18}, F_{19}] = [2584, 4181]\) to ensure that the
learning rate schedule (with Fibonacci-indexed milestones) reaches its
first stable plateau before the first ASHA pruning decision.  The
algebraic connection: \(4000 \in [F_{18}, F_{19}]\) means that the
warmup period spans exactly one Fibonacci generation at the scale relevant
to the training dynamics.

\begin{remark}
This Fibonacci-band argument for the warmup count is heuristic, not
formally proved.  The Coq certificate for the warmup constant is INV-2
(\texttt{prune\_threshold\_from\_trinity} in
\filepath{trinity-clara/proofs/igla/igla\_asha\_bound.v}), which bounds
the pruning threshold but does not directly justify 4000 as a Fibonacci
band member.  We record this as an open algebraic question (R5 honesty).
\end{remark}

% ============================================================
\section{Philosophical Coda: Recurrence, Morphism, and Growth}
\label{sec:ch7-coda}
% ============================================================

The title ``Golden Sprout'' captures a biological metaphor: a sprout is
the first visible expression of a seed's growth rule.  In our setting, the
``seed'' is the equation \(x^2 = x + 1\), the ``growth rule'' is the
Fibonacci recurrence \(F_{n+1} = F_n + F_{n-1}\), and the ``sprout'' is
the ring \(\mathbb{Z}[\varphi]\) that emerges when we take the unique
positive root \(\varphi\) seriously as an algebraic object, not merely a
numerical approximation.

The Rule of Three in this chapter maps onto three aspects of growth:

\begin{enumerate}
  \item \textbf{Additive growth} (Strand~I): each generation is the sum
        of the previous two, realising the principle of ``growth by
        aggregation'' that Fibonacci observed in the rabbit-breeding problem.
  \item \textbf{Multiplicative growth} (Strand~II): each generation is
        the \(\varphi\)-fold scaling of the previous, realising the principle
        of ``growth by ratio'' that the golden spiral embodies.
  \item \textbf{Modular growth} (Strand~III): in finite fields, growth
        wraps around, realising the principle of ``growth under constraint''
        that the GF16 mantissa architecture exploits for bounded arithmetic.
\end{enumerate}

The Trinity anchor \(\phipow{2} + \phipow{-2} = 3\) is the equation that
makes all three simultaneously possible: it is the algebraic signature that
\(\varphi\) is not just a ratio but a ring element, and that the Fibonacci
sequence is not just a combinatorial accident but a morphism orbit.

\begin{quotation}
\emph{``The Fibonacci numbers are not just a sequence; they are the integer
coordinates of a trajectory in an algebraic orbit.''} ---
Paraphrasing the spirit of \cite{hardy_wright}, Chapter~10.
\end{quotation}

This philosophical observation has a precise technical counterpart:
Theorem~\ref{thm:fib-morphism} states that the Fibonacci recurrence is
exactly the ring-endomorphism recurrence of \(\mathbb{Z}[\varphi]\).
Nothing more is claimed.  No metaphysical weight is attached to the
golden ratio beyond its algebraic properties as the positive root of
\(x^2 = x + 1\) (R5 honesty).

% ============================================================
\section{Selected Exercises and Open Problems}
\label{sec:ch7-exercises}
% ============================================================

The following exercises and open problems consolidate the chapter's results
and point toward future work.  They are intended for a graduate reader who
wishes to deepen their understanding of the \(\mathbb{Z}[\varphi]\) structure.

\subsection{Exercises}\label{subsec:ch7-exercises-list}

\begin{enumerate}

\item \textbf{(Companion matrix mod 5)}  Compute \(\mathbf{M}^{20} \pmod 5\)
directly and confirm that it equals the identity matrix modulo 5,
verifying \(\pi(5) = 20\) via Lemma~\ref{lem:matrix-power}.

\item \textbf{(Norm computation)}  Let \(\alpha = 3 + 5\varphi\).  Compute
\(N(\alpha) = 3^2 + 3 \cdot 5 - 5^2\) and determine whether \(\alpha\) is a
unit, an irreducible, or a composite element of \(\mathbb{Z}[\varphi]\).

\item \textbf{(Zeckendorf of 100)}  Write 100 as a Zeckendorf sum of
non-consecutive Fibonacci numbers.  Express the result as an element
\(\alpha \in \mathbb{Z}[\varphi]\) and verify that \(N(\alpha) \in \mathbb{Z}\).

\item \textbf{(Trinity trace)}  Prove directly that
\(\mathrm{Tr}(\phipow{n}) = L_n\) for \(n \in \{0,1,2,3,4\}\)
by computing both sides from the definitions and checking agreement.

\item \textbf{(Morphism of Lucas)}  Show that the map
\(\Lambda : \mathbb{Z}[\varphi] \to \mathbb{Z}[\varphi]\) defined by
\(\Lambda(\alpha) = \alpha + \sigma(\alpha)\) (where \(\sigma\) is Galois
conjugation) is a ring homomorphism from \(\mathbb{Z}[\varphi]\) to its
subring \(\mathbb{Z}\).  Use this to derive the Lucas recurrence from the
Fibonacci recurrence.

\item \textbf{(Pisano periods)}  Compute \(\pi(2)\), \(\pi(3)\), \(\pi(4)\),
\(\pi(5)\), and \(\pi(8)\) by direct enumeration.  Verify that
\(\pi(16) \leq \pi(8) \cdot \pi(2)\) using the formula for Pisano periods
of prime powers.

\item \textbf{(GF16 multiplication table)}  Using Table~\ref{tab:zeckendorf-gf16},
compute \(5 \cdot 8\), \(3 \cdot 7\), and \(11 \cdot 13\) in the
Zeckendorf representation, reducing modulo 16.

\item \textbf{(Cassini in \(\mathbb{Z}[\varphi]\))}  Express Cassini's identity
(Proposition~\ref{prop:cassini}) as the statement that
\(N(\phipow{n}) = (-1)^n\), and deduce it from
Proposition~\ref{prop:units}.

\item \textbf{(Continued fraction of \(\sqrt{5}\))}  Show that
\(\sqrt{5} = [2; 4, 4, 4, \ldots]\) and derive from this that
\(\varphi = [1; 1, 1, 1, \ldots]\).  What is the continued fraction
of \(\psi = (1-\sqrt{5})/2\)?

\item \textbf{(Fibonacci entry point of 89)}  \(89\) is itself a Fibonacci
number (\(F_{11} = 89\)).  Determine the entry point \(\alpha(89)\) and the
Pisano period \(\pi(89)\).  Verify that \(89 \equiv 4 \pmod 5\), so
\(89\) is inert in \(\mathbb{Z}[\varphi]\), and use this to bound
\(\alpha(89)\).

\end{enumerate}

\subsection{Open Problems}\label{subsec:ch7-open}

\begin{enumerate}

\item \textbf{(fib\_tess\_closure, INV-12)}  Formalise
Claim~\ref{claim:fib-tess} in Coq.  The file
\filepath{trinity-clara/proofs/fib\_tess.v} does not yet exist.  A minimal
Coq proof would: (a) import the Lucas closure lemma from
\filepath{lucas\_closure\_gf16.v}; (b) establish the Pisano period of 16
computed symbolically; (c) prove that the companion matrix
\(\mathbf{M}^{24} \equiv I \pmod{16}\); (d) conclude closure.

\item \textbf{(Binet in Coq)}  Formalise Theorem~\ref{thm:binet} within
the \texttt{trinity-clara} proof environment.  The main challenge is
working with \(\sqrt{5}\) in Coq's real number library \texttt{Coq.Reals}.
A possible approach: avoid \(\sqrt{5}\) entirely by working in the ring
\(\mathbb{Z}[x]/(x^2-x-1)\) and proving the formula symbolically, then
specialising to \(x = \varphi\).

\item \textbf{(Zeckendorf in Coq)}  Formalise Theorem~\ref{thm:zeckendorf}
in Coq.  The standard inductive proof is straightforward, but the
uniqueness argument requires a formalised theory of well-founded induction
on Fibonacci indices.  The \texttt{mathcomp} library provides the necessary
tools.

\item \textbf{(Wall--Sun--Sun and GF16)}  Make precise the heuristic
connection between Wall--Sun--Sun primes and GF16 closure failures
(Remark~\ref{rem:wss}).  Specifically: if \(p\) is a Wall--Sun--Sun prime,
does the GF\(p\) mantissa system (the \(p\)-value analogue of GF16) fail
the tesselation closure property?  This would make the non-existence of
Wall--Sun--Sun primes a necessary condition for the universality of the
Golden Sprout quantisation scheme.

\item \textbf{(Lucas ring and GF16 exponents)}  In Chapter~6, the GF16
exponent uses a 3-band structure derived from
\(\phipow{2}+\phipow{-2}=3\).  Formalise the correspondence between
the three exponent bands and the three Lucas-ring generators
\(\{1, \phipow{2}, \phipow{-2}\}\) of the exponent lattice.

\end{enumerate}

% ============================================================
\section{References}\label{sec:ch7-refs}
% ============================================================

\begin{thebibliography}{99}

\bibitem{koshy_fib_lucas}
T.~Koshy,
\emph{Fibonacci and Lucas Numbers with Applications}, 2nd~ed.
Wiley, 2018.
ISBN 978-1118742129.

\bibitem{zeckendorf1972}
E.~Zeckendorf,
``Repr\'esentation des nombres naturels par une somme de nombres de
Fibonacci ou de nombres de Lucas,''
\emph{Bulletin de la Soci\'et\'e Royale des Sciences de Li\`ege},
vol.~41, pp.~179--182, 1972.
CODEN~FOCRDQ.

\bibitem{hardy_wright}
G.~H.~Hardy and E.~M.~Wright,
\emph{An Introduction to the Theory of Numbers}, 6th~ed.
Oxford University Press, 2008.
ISBN 978-0199219865.

\end{thebibliography}
